\clearpage\SourceLeaf{289}
\begin{center}
[II.]
\end{center}

\begin{center}\textbf{ZUR GEOMETRIE DER LAGE,}\par\medskip\textbf{FÜR ZWEI RAUMDIMENSIONEN.}\end{center}

Bei den dargestellten Tracten werden alle nur einmal vorkommenden Knotenpunkte weggelassen. Mit zwei Sternen sind diejenigen bezeichnet, die für sich schon unmöglich sind; mit Einem Stern die, wo unter Vor- und Nachsetzung eines neuen Knotenpunkts der Tract unmöglich; ohne Stern, wo dieser Zusatz einen möglichen Tract ergibt.

\begin{center}
Ein vollständiger Knoten.

1. \emph{aa}
\end{center}

\begin{center}
Zwei vollständige Knoten.

1. \emph{aabb}\\
2. \emph{abab} *\\
3. \emph{abba}
\end{center}

\begin{center}
Drei vollständige Knoten.
\end{center}
\begin{center}
\begin{minipage}[t]{.30\linewidth}\centering\begin{tabular}[t]{rl}
1. & \emph{aabbcc} \\
2. & \emph{aabcbc}\,* \\
3. & \emph{aabccb} \\
4. & \emph{ababcc}\,* \\
5. & \emph{abacbc}\,** \\
6. & \emph{abaccb}\,* \\
7. & \emph{abbacc} \\
8. & \emph{abbcac}\,* \\
\end{tabular}\end{minipage}\hspace{1em}\begin{minipage}[t]{.30\linewidth}\centering\begin{tabular}[t]{rl}
9. & \emph{abbcca} \\
10. & \emph{abcabc} \\
11. & \emph{abcacb}\,* \\
12. & \emph{abcbac}\,* \\
13. & \emph{abcbca}\,** \\
14. & \emph{abccab}\,* \\
15. & \emph{abccba} \\
\end{tabular}\end{minipage}
\end{center}
